\label{ch:conclusions}

This thesis investigated a narrow systems problem created by dynamic expert
placement: how to migrate expert weights without allowing background network
traffic to interfere unnecessarily with ongoing inference. The key conceptual
step was to distinguish migration completion time from foreground serving
quality. In synchronous EPLB, finishing the migration quickly is appropriate
because inference is paused. In asynchronous EPLB, inference and migration
share resources, so lower peak contention can be worth additional transfer
rounds.

An offline ILP formalized this trade-off using flow-to-batch assignments,
rank-degree constraints, routed link load, and a worst-case inference-time
surrogate. It computes a constrained optimum after imposing the same
rank-disjoint-flow policy as the deployed design; it does not establish that
this concurrency limit is universally optimal. In an illustrative hotspot, the
constrained schedule reduced normalized
peak utilization from 1.5 to 0.5 and modeled inference time from 4.75 to 2.25,
while increasing estimated migration time from 1.55 to 2.20. The purpose of the
model was explanatory: it established that the apparent contradiction between
``more batches'' and ``faster serving'' disappears when the objectives are
separated.

The deployed scheduler translates that insight into deterministic first-fit
edge coloring. Remote migrations are coalesced by ordered rank pair, and every
batch forms a matching in the rank communication graph. Independent pairs
retain concurrency, but no rank participates in multiple directed flows within
a batch.
All layer schedules are constructed in one cycle-wide CPU call, while expert
weights remain transferred and activated layer by layer. The implementation is
opt-in, async-only, and preserves the legacy path as a safe fallback.

On the evaluated four-node vLLM system, enabling batching improved output
throughput by 2.28\%, reduced TTFT P50/P99 by 1.18\%/1.30\%, reduced TPOT
P50/P99 by 1.94\%/0.87\%, and reduced E2E P50/P99 by 3.18\%/1.50\%. Head-node
NIC P99 fell by 10.24\%. Profiling measured 34 scheduling calls for all 48
layers per rank, with \SI{302.081}{ms} as the largest cumulative rank cost.
That cost was 0.285\% of the independently observed \SI{105.932}{s} serving-time
reduction. The code was merged upstream into vLLM in pull request \prnumber.

The work also has clear limitations. It reports a single A/B pair on one
cluster, uses a synthetic random workload, samples NIC traffic on only the head
node, and evaluates one migration backend. Accordingly, the result is an
engineering case study rather than a universal throughput claim.

Several extensions follow naturally:

\begin{itemize}
  \item repeat and randomize A/B runs on server-class GPUs and RDMA networks,
  reporting confidence intervals and per-node network counters;
  \item relax the one-flow-per-rank limit using observed foreground load or
  available bandwidth rather than fixing it statically;
  \item incorporate topology, heterogeneous expert sizes, and multiple NICs
  while retaining a deterministic fallback;
  \item evaluate synchronization semantics and benefit across every supported
  communicator backend;
  \item coordinate placement and migration scheduling so a placement policy can
  account for reconfiguration cost without coupling model routing to a complex
  online optimizer; and
  \item expose migration progress and interference metrics for automatic
  operator tuning.
\end{itemize}

The broader lesson is that online model reconfiguration should be evaluated as
part of serving, not as an isolated data-transfer benchmark. A migration that
finishes later can still be the better systems choice when it protects the
latency and throughput of user requests.
